%% ============================================================
%% CHAPTER 8 — DISCUSSION
%% ============================================================
\chapter{Discussion}
\label{ch:discussion}

The four results chapters tell one story. Under consistent, second-law-
consistent assumptions and with the waste-gas stream treated as what it
is --- a fixed resource at a warm site --- the steam Rankine cycle is the
most efficient, the cheapest and the most bankable option at 350 and
500\,\textdegree C, the transcritical CO$_2$ cycle recovers the most energy
and is the most flexible option but needs a cold heat sink that Australia
does not naturally provide, and the supercritical CO$_2$ Brayton cycle, in
its simple recuperated form, is not competitive below 500\,\textdegree C at
all. This chapter explains why these conclusions differ from much of the
published literature and from this project's own Interim Report, sets out
the figure of merit that should be used for WHR cycle selection, develops
the Australian heat-sink constraint, and states the limitations within which
the conclusions hold.

\section{Why the Literature's CO$_2$ Advantage Did Not Appear}
\label{sec:discussion_paradox}

Four boundary conditions account for the difference between this study and
the 5--15 point efficiency advantages reported for CO$_2$ cycles by Kizilkan
\cite{Kizilkan2020}, Jiang et al.~\cite{Jiang2023} and Crespi et
al.~\cite{Crespi2017}; each is quantified in Chapter~\ref{ch:results_rq1}.

\begin{enumerate}[leftmargin=1.8em,itemsep=4pt]
\item \textbf{Recuperator energy balance (Finding~3).} With the effectiveness
applied to the limiting stream, recuperation adds 3--10 points to the
tCO$_2$ cycle, not 12--20. Near the critical point the high-pressure stream
can absorb two to three times more enthalpy per kelvin than the low-pressure
stream can supply; a cold-side-only effectiveness transfers heat the exhaust
does not have. The Interim Report made this error, and some published
models that report unusually high recuperated efficiencies may as well.
\item \textbf{Heat sink temperature (Findings~11--13).} The sCO$_2$ cycle's
efficiency rests on compressing a dense near-critical fluid; at a
36\,\textdegree C compressor inlet, which is what Australian cooling water
allows, the fluid is already gas-like and the back-work ratio is 31--64\,\%
at the three sites. Moving the inlet from 32 to 40\,\textdegree C removes
two-fifths of the cycle's efficiency at 200\,\textdegree C. The tCO$_2$
baseline condenses at 14\,\textdegree C, which no Australian site can
supply without refrigeration.
\item \textbf{Steam baseline.} Many comparisons evaluate steam at a single
modest pressure. At turbine inlet temperatures of 350--500\,\textdegree C a
10\,MPa steam cycle with a 10\,kPa condenser reaches 32--34\,\% and is a
demanding baseline (Finding~1). When the boiler is integrated with the
stream the steam pressure falls to 0.9--3.7\,MPa and the site efficiency to
24--31\,\% (Finding~14), close to Kizilkan's 24\,\%; the CO$_2$ cycles then
compete on recovery depth, not on efficiency.
\item \textbf{Cycle layout.} The simple recuperated sCO$_2$ layout
modelled here is the standard single-loop configuration. Its recuperated
heater inlet of 100--279\,\textdegree C means it recovers only 60--65\,\% of
the heat that the tCO$_2$ cycle recovers from the same stream (Finding~15). Cascade, dual-recuperated and preheating layouts
that address this were outside the scope of a consistent three-cycle
comparison and are the obvious next step (Section~\ref{sec:future_work}).
\end{enumerate}

The disagreement, then, is not with the thermodynamics of CO$_2$ but with
the conditions under which it is attractive: a cold sink, a high source
temperature, and a layout designed for deep recovery. Australian industrial
WHR at 200--500\,\textdegree C offers none of the first two.

\section{The Figure of Merit for Waste Heat Recovery}
\label{sec:discussion_fom}

Thermal efficiency is the wrong single number for WHR cycle selection. The
owner of a 50\,kg/s cement exhaust does not care what fraction of the
heat entering the cycle becomes electricity; they care how much electricity
the exhaust yields and what the plant costs. Equation~\eqref{eq:wgas},
$w_{gas} = \eta_{th}\,c_{p,gas}(T_{src} - T_{gas,out})$, captures the first
quantity, and the results of Chapters~\ref{ch:results_rq2}
and~\ref{ch:results_rq3} show that the second factor can outweigh the
first: tCO$_2$ at CS2 is 1.2~points more efficient than steam but recovers
75\,\% more power; sCO$_2$ at CS1 is 1.4~points less efficient than steam but
recovers 25\,\% less.

The second quantity, cost, must then be evaluated at the output the cycle
actually delivers, because specific capital cost falls steeply with size
(Table~\ref{tab:capital}). A fixed-output comparison at 1\,MWe --- the basis
of the Interim Report and of much of the literature --- reduces the whole
economic analysis to the capital cost ranking, because with free heat and a
common life and capacity factor LCOE is proportional to TEC/kW
(Finding~25). The fixed-source framing used here is the one a feasibility
study would use, and it is the one on which tCO$_2$ earns its single
economic win, at the 200\,\textdegree C site.

\section{Heat Rejection in a Warm Climate}
\label{sec:discussion_australian}

The decisive Australian constraint is the heat sink. Three results combine:

\begin{itemize}[leftmargin=1.5em,itemsep=3pt]
\item The tCO$_2$ cycle at $P_{low} = 5$\,MPa condenses at 14.3\,\textdegree C
and needs a sink near 4\,\textdegree C (Finding~11). Cooling-tower water at
Gladstone, Sydney or the Pilbara is 20--32\,\textdegree C for most of the
year~\cite{BOM2024}; seawater at the Pilbara coast exceeds 25\,\textdegree C.
\item Supplying that sink with an electric chiller would consume 63--134\,\%
of the cycle's net output at COP~4 (Finding~12). The tCO$_2$ NPV advantage
at CS2 (Finding~23) therefore cannot be realised there; it exists only where
a natural cold sink is available --- cold river or groundwater, a deep-water
intake, or a co-located refrigeration load.
\item Raising $P_{low}$ to 7\,MPa (condensing at 28.7\,\textdegree C) makes
the cycle compatible with evaporative cooling outside summer but costs
23--30\,\% of the recovered work; at CS1 the result is exactly the steam
cycle's 52.8\,kJ per kilogram of gas, at a higher capital cost. At the
Pilbara, with summer wet-bulb temperatures above 28\,\textdegree C, even
7\,MPa is marginal.
\end{itemize}

The exergy analysis flags the same issue from the other direction: the
negative condenser term in Table~\ref{tab:exergy_sites} is the unpaid
refrigeration exergy. The sCO$_2$ cycle faces a milder version of the
constraint at its compressor inlet (Finding~13). The steam cycle, condensing
at 45.8\,\textdegree C, is the only one of the three that is indifferent to
Australian ambient conditions, and this, more than any thermodynamic
property, is why it emerges as the default choice.

\section{Cycle Selection Framework for Australian Industry}
\label{sec:discussion_framework}

Table~\ref{tab:selection} condenses the thirty findings into a selection
guide for a 0.3--3.5\,MWe plant at 2026 costs and prices.

\begin{table}[H]
\centering
\caption{Recommended cycle selection for Australian industrial WHR
(2026 conditions, fixed-source basis).}
\label{tab:selection}
\renewcommand{\arraystretch}{1.3}
\small
\begin{tabularx}{\textwidth}{>{\raggedright\arraybackslash}p{2.6cm}
  >{\raggedright\arraybackslash}p{2.4cm} >{\raggedright\arraybackslash}X}
\toprule
\rowcolor{tableheader}
\textbf{Source} & \textbf{Recommend} & \textbf{Basis and conditions} \\
\midrule
\rowcolor{rowgray}
$\ge$\,450\,\textdegree C (EAF, kiln) & SRC &
LCOE AUD\,31/MWh, NPV 12.4\,M, payback 3.2\,yr, 100\,\% Monte Carlo
viability; steam is also the most efficient cycle here. tCO$_2$ only if load
is highly variable and a cold sink exists. \\
300--450\,\textdegree C (cement) & SRC &
LCOE AUD\,35/MWh, NPV 11.4\,M, payback 4.2\,yr, 99.8\,\% viability. tCO$_2$
recovers 29\,\% more power but costs 94\,\% more; its NPV is 12\,\% lower and
its sink requirement is unmet. \\
\rowcolor{rowgray}
150--300\,\textdegree C (gas processing) & SRC, or tCO$_2$ with a natural cold sink &
Both marginal: LCOE 96--100 against a 130\,AUD/MWh price, payback 10--11\,yr,
68--72\,\% viability. tCO$_2$ gives 53\,\% more NPV where water below
$\approx$10\,\textdegree C is available; otherwise steam, or no project. \\
Any temperature & Not sCO$_2$ (simple recuperated) &
Negative or marginal NPV at every site; $P(\mathrm{NPV}>0)$ 0.5--49\,\%.
Reconsider only with a cascade layout, a compressor inlet below
33\,\textdegree C, and sources above 500\,\textdegree C. \\
\rowcolor{rowgray}
Fluctuating load & tCO$_2$ preferred & Marchionni score 0.78 against 0.73
(SRC) and 0.71 (sCO$_2$); steam is weakest at 200\,\textdegree C. \\
\bottomrule
\end{tabularx}
\end{table}

Three qualifications apply. First, the steam recommendation at
200\,\textdegree C is weak: with an IRR of 12\,\% against a 7\,\% discount
rate and a 28\,\% chance of loss, the project depends on the Pilbara price
premium and would not proceed at NEM prices. Second, the tCO$_2$ option at
any site should be evaluated with the actual sink: with cold water the cycle
is the best choice at 200\,\textdegree C; with a chiller it is not viable;
at 7\,MPa it is a more expensive equivalent of steam. Third, all three
cycles are small plants in cost terms --- heat exchangers are 35--43\,\% of
CO$_2$ equipment cost and heat rejection 17--44\,\% of every cycle's --- so
the recommendation is sensitive to the heat exchanger correlations, which is
why the Monte Carlo range of $\pm30$\,\% on capital was carried through.

\section{Limitations and Assumptions}
\label{sec:discussion_limitations}

\begin{limitbox}
\textbf{L1 --- Constant turbomachinery efficiencies.} $\eta_t = 0.85$ and
$\eta_{p/c} = 0.80$ are applied at all sizes and all three cycles. Real
sub-MW machines are less efficient, and the penalty is probably larger for
the high-pressure CO$_2$ turbines than for steam.

\textbf{L2 --- No pressure drop.} All exchangers are isobaric. PCHE pressure
drops of 1--2\,\% of the high-side pressure would reduce the CO$_2$ cycle
efficiencies by roughly one point and raise the steam cycle's by nothing.

\textbf{L3 --- Cost correlations extrapolated.} The Weiland turbine and
generator correlations begin at 8 and 4\,MW; the Turton compressor at
450\,kW. Their stated uncertainties ($-32$ to $+51$\,\% on the turbine) are
wider than the $\pm30$\,\% carried in the Monte Carlo. The steam turbine
benchmark is for back-pressure units; a condensing turbine with VFR of
45--148 would cost more, which would narrow but not close the gap to tCO$_2$
at CS1 and CS3.

\textbf{L4 --- Heat rejection modelled uniformly.} A 10\,K mean temperature
difference and the Weiland air-cooler correlation are applied to all three
cycles. This is conservative for the 46\,\textdegree C steam condenser
(cheaper water-cooled units exist) and optimistic for the 14\,\textdegree C
tCO$_2$ condenser (no chiller is costed; Finding~12 quantifies the omission).

\textbf{L5 --- CO$_2$ pressures not re-optimised for recovery.} The CO$_2$
high- and low-side pressures follow the RQ1 rules, whereas the steam
pressure was optimised for $w_{gas}$ at each site. A tCO$_2$ high-pressure
optimisation on the same basis would improve its results by a few per cent.

\textbf{L6 --- Simple cycle layouts.} One layout per cycle, no reheat,
recompression, cascade or split flow. This is the right basis for a
like-for-like comparison but understates what an optimised sCO$_2$ system
could do above 500\,\textdegree C.

\textbf{L7 --- Monte Carlo scope.} Capital, price and capacity factor are
varied; source temperature and flow, discount rate, O\&M and plant life are
not. Correlation between variables is ignored.

\textbf{L8 --- Part-load model.} A single parabolic Stodola coefficient
calibrated to one tCO$_2$ dataset~\cite{Yang2024} is applied to all three
cycles, and the heater is not re-integrated at part load.

\textbf{L9 --- Economic scope.} No carbon price, no grid-connection or
site-integration costs, no availability penalty for the immature CO$_2$
technologies, and an 85\,\% capacity factor that may be optimistic for
batch processes.
\end{limitbox}

None of these limitations favours the CO$_2$ cycles systematically except
L4 (no chiller) and L3 (back-pressure steam benchmark), which act in opposite
directions; L1, L2 and L3 (turbine extrapolation) act against them. The
qualitative conclusions --- steam first above 300\,\textdegree C, tCO$_2$
viable only with a cold sink, simple sCO$_2$ not viable --- are unlikely to
change within these uncertainties; the margins would.

\section{Comparison With Published Literature}
\label{sec:discussion_comparison}

The present results agree with Baglietto et al.~\cite{Baglietto2023} that the
mature steam supply chain gives steam a capital cost advantage that CO$_2$
cycles do not overcome at industrial WHR scale, and extend that finding to
the smaller 0.3--3.5\,MWe range and to Australian prices. They agree with
Kizilkan~\cite{Kizilkan2020} on the steam efficiency at a 350\,\textdegree C
cement source once the boiler is integrated with the stream, but not on the
sCO$_2$ advantage, for the reasons given in
Section~\ref{sec:discussion_paradox}. They agree in direction with Jiang et
al.~\cite{Jiang2023} that tCO$_2$ outperforms sCO$_2$ at low temperature and
with Yang et al.~\cite{Yang2024} and Marchionni et al.~\cite{Marchionni2020}
that pumped CO$_2$ cycles have the better part-load behaviour. The
observation that recuperated sCO$_2$ cycles recover a waste-heat stream
shallowly is well known in the sCO$_2$ community and is the reason for the
cascade and split-flow WHR layouts described by Crespi et
al.~\cite{Crespi2017}; this thesis quantifies the penalty for the simple
layout against two alternatives on one basis. The heater-integration result
(Findings~6 and~18) is, to the author's knowledge, the first explicit
demonstration that the pseudo-critical $c_p$ peak produces a double-pinched,
mid-bulging profile rather than an interior pinch, and that two-point LMTD
sizing overestimates the tCO$_2$ heater by a factor of two to three at
$\varepsilon = 0$.

\section{Revisions Since the Interim Report}
\label{sec:discussion_revisions}

For transparency, the five corrections made between the Interim Report and
this thesis, and their consequences, are listed in
Table~\ref{tab:revisions}. Each was found by re-checking the model code
against its own outputs after the Interim Report was marked.

\begin{table}[H]
\centering
\caption{Corrections made after the Interim Report.}
\label{tab:revisions}
\renewcommand{\arraystretch}{1.3}
\small
\begin{tabularx}{\textwidth}{>{\raggedright\arraybackslash}p{3.4cm}
  >{\raggedright\arraybackslash}X >{\raggedright\arraybackslash}p{4.6cm}}
\toprule
\rowcolor{tableheader}
\textbf{Correction} & \textbf{What was wrong} & \textbf{Consequence} \\
\midrule
\rowcolor{rowgray}
Recuperator energy balance (Eq.~\ref{eq:recup}) &
Effectiveness applied to the cold stream only; hot stream cooled below the
cold inlet (second-law violation) at every temperature for both CO$_2$ cycles. &
CO$_2$ efficiencies fall 3--11 points; sCO$_2$/SRC crossover at
362\,\textdegree C and ``tCO$_2$ beats SRC everywhere'' withdrawn (Finding~1). \\
Heater gas profile (Eq.~\ref{eq:gas_profile}) &
Gas temperature linear in node index, not in heat duty, in the pinch
diagrams; gave a constant 10\,K gap and a spurious ``mid-exchanger pinch at
129\,\textdegree C''. &
Double pinch at both terminals with a mid-exchanger bulge (Finding~6).
LMTD comparison (Finding~18) unaffected, as it already used the correct
profile. \\
\rowcolor{rowgray}
Heater sizing in the cost model &
sCO$_2$ heater driving force computed from an expression that cancelled to
$-10$\,K (giving 3,400--7,000\,K); steam boiler driving force from an ad hoc
terminal estimate; tCO$_2$ heater at the worst-case 10\,K. &
sCO$_2$ no longer the cheapest cycle; steam cheapest at every site
(Finding~19). \\
Cost basis &
Turbine coefficients not traceable to the cited source; condenser omitted;
compressor costed as a pump. &
Consistent correlations (Table~\ref{tab:cost_basis}); heat rejection
17--44\,\% of equipment cost (Finding~21). \\
\rowcolor{rowgray}
Economic basis &
Fixed 1\,MWe output regardless of stream; simple payback; steam at
0.95\,$P_{sat}$ at a 200\,\textdegree C source (saturation 2\,K below the
source). &
Fixed-source basis with pinch-limited steam pressure; discounted payback;
tCO$_2$ NPV lead at 200\,\textdegree C (Finding~23). \\
\bottomrule
\end{tabularx}
\end{table}
